\section{The exact comparison with earlier layered constructions}
\label{sec:layered-presence-comparison}

The bounded prior-work search found a precise comparison, not an established identification of the full typed programme with earlier work. It did not establish historical novelty. The following calculation uses the carrier and operation formulas from
Zur Izhakian, Manfred Knebusch and Louis Rowen,
\href{https://arxiv.org/abs/0912.1398v3}{arXiv:0912.1398v3},
Construction \texttt{defn5} (original source lines 1325--1341), with the one-versus-two calculation in Proposition \texttt{cancel1} (lines 1366--1374). The root task read these original TeX passages; this independent calculation uses their full operation formulas.  We apply those formulas to the explicitly
specified auxiliary monoid below and check every required identity.
The letters \(t,u\) are auxiliary monoid elements, not names for the
user's typed objects.  The algebra \(C\) considered here is the
positive-count restriction of the integer extension chosen in the FMC
construction.  This comparison neither asserts that this extension is
uniquely forced by the original rules nor assigns a parity to \(\tau\).

Let \(\mathbb N_{>0}=\{1,2,3,\ldots\}\).  Give \(G=\{t,u\}\) the order
\(t<u\) and multiplication
\[
tt=t,\qquad tu=ut=u,\qquad uu=u.
\tag{LPC1}\label{eq:LPC1}
\]
This multiplication is associative and commutative, with identity \(t\):
a product is \(t\) exactly when every factor is \(t\), and is \(u\)
otherwise.  It preserves the order in each argument.  It is
noncancellative, since \(tu=uu\) while \(t\ne u\).

Define the carrier \(H=(\mathbb N_{>0}\times G)\sqcup\{0\}\).
The element \(0\) is an additive identity and a multiplicative absorber.
For positive integers \(k,l\), the two operations on the other elements
are
\[
(k,a)(l,b)=(kl,ab),\qquad
(k,a)+(l,b)=
\begin{cases}
(k,a),&a>b,\\
(l,b),&a<b,\\
(k+l,a),&a=b.
\end{cases}
\tag{LPC2}\label{eq:LPC2}
\]
These are definitions of a two-operation algebra.  Its multiplication is
associative and commutative by the corresponding facts for positive
integers and for \(G\), including the prescribed rule for \(0\).
Its multiplicative identity is \((1,t)\).  Its addition is commutative.
To prove associativity of addition, consider three nonzero elements
\((k_i,a_i)\), \(1\leq i\leq3\), and put
\(a=\max\{a_1,a_2,a_3\}\).  Both parenthesizations give
\((\sum_{i:a_i=a}k_i,a)\): if the first two labels are both below \(a\),
the third term is retained; if exactly one of the first two has label
\(a\), that term is retained before the third is added; and if both have
label \(a\), their counts are added before the third is added.
These cases also cover the other parenthesization after interchanging
the first and third terms.  If a term is \(0\), the additive-identity
rule reduces the assertion to two terms.  This proves associativity in
every case.

Define the disjoint union
\(C=\{0,\tau\}\sqcup\mathbb N_{>0}\).  For positive integers \(n,m\),
use their ordinary integer sum and product, and set
\[
\begin{gathered}
0+x=x+0=x,\qquad 0x=x0=0\quad(x\in C),\\
\tau+\tau=\tau,\qquad \tau+n=n+\tau=n,\\
\tau\tau=\tau,\qquad \tau n=n\tau=n.
\end{gathered}
\tag{LPC3}\label{eq:LPC3}
\]
The parity label is attached to each positive integer \(n\) by its
integer parity.  In particular, \(\tau\) has no parity label.

Both operations on \(C\) are commutative.  For associativity of addition,
if at least one of three terms is a positive integer, either
parenthesization gives the ordinary sum of precisely the positive
integer terms; \(0\) and \(\tau\) contribute according to the displayed
rules.  If there is no positive integer term but there is a \(\tau\),
both parenthesizations give \(\tau\).  If all terms are zero, both give
zero.  These cases exhaust all triples.  For associativity of
multiplication, a zero factor makes both parenthesizations zero.  Among
nonzero factors, if there is a positive integer factor, both
parenthesizations give the ordinary product of precisely the positive
integer factors; if all factors are \(\tau\), both give \(\tau\).
Thus multiplication is also associative.  Its identity is \(\tau\).

Define maps of carriers \(Q:H\longrightarrow C\) and
\(S:C\longrightarrow H\) by
\[
\begin{gathered}
Q(0)=0,\qquad Q(k,t)=\tau,\qquad Q(k,u)=k,\\
S(0)=0,\qquad S(\tau)=(1,t),\qquad S(n)=(n,u).
\end{gathered}
\tag{LPC4}\label{eq:LPC4}
\]
Every value is in the indicated carrier.  Direct substitution gives
\[
\begin{gathered}
Q\circ S=\operatorname{id}_C,\qquad
(S\circ Q)(0)=0,\\
(S\circ Q)(k,t)=(1,t),\qquad
(S\circ Q)(k,u)=(k,u).
\end{gathered}
\tag{LPC5}\label{eq:LPC5}
\]
Thus \(Q\) is surjective, \(S\) is injective, and \(S\circ Q\) is
idempotent.  All fibres of these maps are
\[
\begin{gathered}
Q^{-1}(\{0\})=\{0\},\qquad
Q^{-1}(\{\tau\})=\{(k,t):k\in\mathbb N_{>0}\},\qquad
Q^{-1}(\{n\})=\{(n,u)\},\\
S^{-1}(\{0\})=\{0\},\qquad
S^{-1}(\{(1,t)\})=\{\tau\},\\
S^{-1}(\{(k,t)\})=\varnothing\quad(k\geq2),\qquad
S^{-1}(\{(n,u)\})=\{n\}.
\end{gathered}
\tag{LPC6}\label{eq:LPC6}
\]

The map \(Q\) preserves addition in every case.  For nonzero terms the
complete calculation, up to interchanging the terms, is
\[
\begin{aligned}
Q\bigl((k,t)+(l,t)\bigr)&=Q(k+l,t)=\tau=\tau+\tau,\\
Q\bigl((k,u)+(l,u)\bigr)&=Q(k+l,u)=k+l,\\
Q\bigl((k,t)+(l,u)\bigr)&=Q(l,u)=l=\tau+l.
\end{aligned}
\tag{LPC7}\label{eq:LPC7}
\]
An input zero gives the same assertion by the additive-identity rules.

The exact comparison of products is
\[
\begin{array}{c|c|c}
(x,y)&Q(xy)&Q(x)Q(y)\\ \hline
((k,t),(l,t))&\tau&\tau\\
((k,u),(l,u))&kl&kl\\
((k,t),(l,u))&kl&l\\
((k,u),(l,t))&kl&k\\
(0,y)\text{ or }(x,0)&0&0
\end{array}
\tag{LPC8}\label{eq:LPC8}
\]
The third row agrees exactly when \(k=1\), since \(l\) is a positive
integer.  The fourth agrees exactly when \(l=1\), since \(k\) is a
positive integer.  Hence all failures of multiplicativity are precisely
the mixed-label products whose \(t\)-labelled factor has count at least
two.  The product of two \(t\)-labelled factors is preserved.

The map \(S\) preserves multiplication in every case:
\[
\begin{aligned}
S(\tau\tau)&=(1,t)=(1,t)(1,t)=S(\tau)S(\tau),\\
S(\tau n)&=(n,u)=(1,t)(n,u)=S(\tau)S(n),\\
S(nm)&=(nm,u)=(n,u)(m,u)=S(n)S(m).
\end{aligned}
\tag{LPC9}\label{eq:LPC9}
\]
Commutativity covers the reversed mixed product, and a zero factor is
covered by the absorber rule.  Its additive comparison is
\[
\begin{aligned}
S(\tau+\tau)&=(1,t),&
S(\tau)+S(\tau)&=(2,t),\\
S(\tau+n)&=(n,u),&
S(\tau)+S(n)&=(n,u),\\
S(n+m)&=(n+m,u),&
S(n)+S(m)&=(n+m,u).
\end{aligned}
\tag{LPC10}\label{eq:LPC10}
\]
An input zero is again preserved.  Since \((1,t)\ne(2,t)\), the ordered
pair \((\tau,\tau)\) is the only failure of additivity of \(S\).

There is an explicit failure of distributivity in \(H\):
\[
\begin{aligned}
(1,u)\bigl((1,t)+(1,u)\bigr)
  &=(1,u)(1,u)=(1,u),\\
(1,u)(1,t)+(1,u)(1,u)
  &=(1,u)+(1,u)=(2,u).
\end{aligned}
\tag{LPC11}\label{eq:LPC11}
\]
The distinct labels \(t<u\) become equal after multiplication by \(u\),
because \(ut=uu=u\); the two resulting counts on the second line are
therefore added.  This is the exact cancellation failure responsible
for the displayed discrepancy.  There is also an explicit failure of
distributivity in \(C\):
\[
1(\tau+1)=1\cdot1=1,
\qquad
1\tau+1\cdot1=1+1=2.
\tag{LPC12}\label{eq:LPC12}
\]
Both calculations follow directly from the specified operations.
Consequently neither \(H\) nor \(C\) is a semiring, while each still has
the associative and commutative operations proved above.  No
distributive law was used to define the user's typed rules.

Finally, the fibres of \(Q\) form an additive congruence, since
\(Q(x+y)=Q(x)+Q(y)\).  They do not form a multiplicative congruence:
for every \(k\geq2\),
\[
\begin{gathered}
Q(k,t)=Q(1,t)=\tau,\\
Q\bigl((k,t)(1,u)\bigr)=k\ne1
 =Q\bigl((1,t)(1,u)\bigr).
\end{gathered}
\tag{LPC13}\label{eq:LPC13}
\]
Thus this identification of all \(t\)-labelled counts cannot give a
quotient preserving both operations of \(H\).  Its exact surviving
relationships are the additive map \(Q\), the multiplicative map \(S\),
the identity \(Q\circ S=\operatorname{id}_C\), and the complete failures
recorded in \eqref{eq:LPC8} and \eqref{eq:LPC10}.
